\documentclass[11pt]{article}
\usepackage[margin=1in]{geometry}
\usepackage{amsmath,amssymb,amsthm}
\usepackage{graphicx}
\usepackage[hidelinks]{hyperref}
\setlength{\emergencystretch}{2em}

\newtheorem{theorem}{Theorem}
\newtheorem{lemma}{Lemma}
\newtheorem{definition}{Definition}
\newtheorem{proposition}{Proposition}

\newcommand{\C}{\mathbb C}
\newcommand{\Kcal}{\mathcal K}
\newcommand{\Prim}{\operatorname{Prim}}
\newcommand{\Mcal}{\mathcal M}

\title{FDI Fell Algebras Have the Expected BCP Criterion}
\author{Run \texttt{fa\_banach\_001}, lane 19}
\date{June 17, 2026}

\begin{document}
\maketitle

\section*{Source Problem}

The source paper is Jinghao Huang, Karimbergen Kudaybergenov, and Rui Liu,
\emph{The ball-covering property of von Neumann algebras and noncommutative
symmetric spaces}, arXiv:2311.14257.  After proving that a commutative unital
\(C^*\)-algebra \(C(K)\) has BCP/UBCP exactly when \(K\) has a countable
\(\pi\)-basis, the authors ask for a noncommutative version:

\begin{center}
\includegraphics[width=0.92\textwidth]{figures/open_problem_crop.png}
\end{center}

\section*{Definitions}

\begin{definition}
A \(C^*\)-algebra \(A\) has the ball-covering property, or BCP, if the unit
sphere of \(A\) is covered by countably many open balls none of which contains
the origin.  It has the uniform ball-covering property, or UBCP, if the balls
can be chosen with a common positive gap from the origin.
\end{definition}

\begin{definition}
A \(C^*\)-algebra is \emph{FDI} if all its irreducible representations are
finite dimensional.  Equivalently, every primitive quotient is a full matrix
algebra \(M_n(\C)\), with \(n\) allowed to depend on the primitive ideal.
\end{definition}

\section*{Statement}

\begin{theorem}[FDI Fell criterion]
\label{thm:main}
Let \(A\) be a Fell \(C^*\)-algebra all of whose irreducible representations
are finite dimensional.  Then
\[
        A\text{ has BCP}
        \quad\Longleftrightarrow\quad
        A\text{ has UBCP}
        \quad\Longleftrightarrow\quad
        \Prim(A)\text{ has a countable }\pi\text{-basis}.
\]
\end{theorem}

\section*{Idea of the Proof}

The earlier local compact-chart packet proves UBCP for any \(C^*\)-algebra
whose primitive space has a countable \(\pi\)-basis of open charts on which
every ambient element is a norm-continuous compact-operator field.  The
multiplier obstruction for general type-I algebras is that an element may
restrict to a continuous-trace ideal only as a multiplier field, and multiplier
fields over \(\Kcal(H)\) need not be norm-continuous.

For FDI Fell algebras this obstruction disappears.  Fell spectra are locally
Hausdorff.  On a Hausdorff open patch, the corresponding ideal is Fell with
Hausdorff spectrum, hence continuous trace.  Since the irreducible
representations are finite dimensional, the local fibres are matrix algebras
\(M_n(\C)\).  For matrix fibres, the multiplier algebra is just bounded
norm-continuous matrix fields.  Thus FDI Fell algebras satisfy the compact
chart hypothesis, and the local chart theorem applies.

\section*{Imported Criteria and Standard Facts}

We use the following criterion from the earlier lane-19 local compact-chart
packet.

\begin{proposition}[local compact-chart criterion]
\label{prop:chart}
Let \(B\) be a \(C^*\)-algebra.  Suppose \(\Prim(B)\) has a countable
\(\pi\)-basis \((O_n)\) such that for each \(n\):
\begin{enumerate}
    \item the ideal \(B_{O_n}\) is isomorphic to
          \(C_0(O_n,\Kcal(H_n))\) for a separable nonzero Hilbert space
          \(H_n\);
    \item for every \(b\in B\), the image of \(b\) on \(O_n\) is represented
          by a bounded norm-continuous function
          \(O_n\to \Kcal(H_n)\), recovering the primitive quotient norms.
\end{enumerate}
Then \(B\) has UBCP.
\end{proposition}

We also use two standard Fell-algebra facts.  First, the spectrum of a Fell
algebra is locally Hausdorff.  Second, a Fell algebra with Hausdorff spectrum
has continuous trace.  The first appears, for example, as Lemma 3.2 of
Clark--an Huef--Raeburn \cite{CHR}, and the second is the classical
Dixmier criterion recalled in their introduction.

\section*{Proof}

\begin{lemma}[FDI Fell algebras have compact-norming charts]
\label{lem:charts}
Let \(A\) be an FDI Fell \(C^*\)-algebra.  If \(\Prim(A)\) has a countable
\(\pi\)-basis, then \(A\) satisfies the hypotheses of
Proposition~\ref{prop:chart}.
\end{lemma}

\begin{proof}
Write \(X=\Prim(A)\), and let \(\mathcal P=\{P_m:m\ge1\}\) be a countable
\(\pi\)-basis for \(X\).

We first refine \(\mathcal P\) to Hausdorff continuous-trace trivialization
patches.  Let \(U\subset X\) be nonempty open.  Choose \(x\in U\).  Since
\(A\) is Fell, \(X\) is locally Hausdorff, so there is an open Hausdorff
neighbourhood \(V\) of \(x\) with \(V\subset U\) after shrinking inside
\(U\).  The ideal \(A_V\) is again Fell, and its spectrum is the Hausdorff
space \(V\).  Hence \(A_V\) has continuous trace.

By local triviality of continuous-trace algebras, there is a nonempty open
set \(W\subset V\) such that
\[
        A_W\cong C_0(W,\Kcal(H))
\]
for some Hilbert space \(H\).  Because \(A\) is FDI, every irreducible
representation over \(W\) is finite dimensional; therefore \(H\) is
finite dimensional on this trivialization.  Since \(\mathcal P\) is a
\(\pi\)-basis for \(X\), there is \(P_m\in\mathcal P\) with
\[
        \emptyset\ne P_m\subset W\subset U.
\]
Restricting the above trivialization gives
\[
        A_{P_m}\cong C_0(P_m,M_d(\C))
\]
for \(d=\dim H<\infty\).  Thus the subfamily of all \(P_m\)'s that occur in
such finite-dimensional continuous-trace trivializations is still a countable
\(\pi\)-basis of \(X\).

Fix one such chart \(O=P_m\), and identify
\[
        A_O\cong C_0(O,M_d(\C)).
\]
Let \(a\in A\).  The canonical action of \(A\) on the ideal \(A_O\) gives a
multiplier
\[
        \phi_O(a)\in \Mcal(A_O).
\]
Since \(M_d(\C)\) is unital and finite dimensional,
\[
        \Mcal(C_0(O,M_d(\C)))\cong C_b(O,M_d(\C)),
\]
with the usual norm-continuous bounded sections.  Let
\[
        \widehat a_O:O\to M_d(\C)
\]
be the section corresponding to \(\phi_O(a)\).  This section is bounded and
norm-continuous.

For \(P\in O\), evaluation at \(P\) gives the irreducible representation of
the ideal \(A_O\), and the multiplier action is exactly the restriction of
the irreducible representation of \(A\) with kernel \(P\).  Hence
\[
        \|a+P\|=\|\widehat a_O(P)\|.
\]
Since \(M_d(\C)=\Kcal(\C^d)\), the chart is a compact-norming chart in the
sense of Proposition~\ref{prop:chart}.  The countable family constructed
above therefore verifies all hypotheses of Proposition~\ref{prop:chart}.
\end{proof}

\begin{proof}[Proof of Theorem~\ref{thm:main}]
If \(\Prim(A)\) has a countable \(\pi\)-basis, then
Lemma~\ref{lem:charts} verifies the compact-norming chart hypothesis for
\(A\).  Proposition~\ref{prop:chart} gives UBCP.

The implication UBCP \(\Rightarrow\) BCP is immediate.  Conversely, if \(A\)
has BCP, the primitive \(\pi\)-base obstruction from the earlier lane-19
packet gives a countable \(\pi\)-basis for \(\Prim(A)\).  This proves the
equivalence.
\end{proof}

\section*{Scope and Limitations}

This theorem is a positive non-Hausdorff Fell/type-I result, but it is not the
full Fell theorem.  The finite-dimensionality of irreducible representations
is doing real work: it makes
\[
        \Mcal(C_0(O,M_d(\C)))=C_b(O,M_d(\C))
\]
a norm-continuous matrix-field algebra.  For infinite-dimensional compact
fibres \(C_0(O,\Kcal(H))\), multiplier fields live in a strict-continuous
\(B(H)\)-bundle, and pointwise compactness or Fell local rank-one behaviour
does not automatically give norm-continuous compact sections.  That
strict-versus-norm gap is the remaining obstacle for Fell algebras with
infinite-dimensional separable irreducible representations.

\section*{Novelty and Search Bounds}

The local indexes were searched for arXiv:2311.14257, FDI, subhomogeneous
Fell algebras, finite-dimensional irreducible representations, local compact
charts, BCP, and UBCP.  Web searches on June 17, 2026 for Fell algebras,
locally Hausdorff spectrum, continuous trace patches, and arXiv:1207.2540
located the Clark--an Huef--Raeburn Fell reference \cite{CHR}, but no BCP
criterion for FDI Fell algebras.  The result is therefore recorded as an
agent-produced partial theorem, subject to human literature review.

\section*{Human Review Focus}

The proof is short because it delegates the covering construction to the
local compact-chart packet.  The main review points are:
\begin{enumerate}
    \item the refinement of a countable \(\pi\)-basis to Hausdorff
          continuous-trace trivialization patches;
    \item the claim that finite-dimensional fibres turn local multiplier
          sections into norm-continuous matrix-valued functions;
    \item the norm-recovery identity
          \(\|a+P\|=\|\widehat a_O(P)\|\) for ambient elements acting as
          multipliers of the local ideal.
\end{enumerate}

\begin{thebibliography}{9}

\bibitem{HKL}
Jinghao Huang, Karimbergen Kudaybergenov, and Rui Liu,
\emph{The ball-covering property of von Neumann algebras and noncommutative
symmetric spaces},
arXiv:2311.14257, 2023.

\bibitem{CHR}
Lisa Orloff Clark, Astrid an Huef, and Iain Raeburn,
\emph{The equivalence relations of local homeomorphisms and Fell algebras},
arXiv:1207.2540, 2012.

\bibitem{Dixmier}
J. Dixmier,
\emph{\(C^*\)-Algebras},
North-Holland, 1977.

\bibitem{RW}
I. Raeburn and D. P. Williams,
\emph{Morita Equivalence and Continuous-Trace \(C^*\)-Algebras},
American Mathematical Society, 1998.

\bibitem{lane19chart}
Run \texttt{fa\_banach\_001}, lane 19,
\emph{A Local Compact-Chart Criterion for \(C^*\)-Algebras with UBCP},
solution packet
\texttt{2311.14257\_local\_compact\_chart\_bcp\_criterion}, 2026.

\bibitem{lane19primitive}
Run \texttt{fa\_banach\_001}, lane 19,
\emph{A Primitive-Space \(\pi\)-Base Obstruction for \(C^*\)-Algebras with
BCP},
solution packet
\texttt{2311.14257\_primitive\_pi\_base\_obstruction}, 2026.

\end{thebibliography}

\end{document}
